\subsection{An obstruction with infinite nonunit coefficient support}

\begin{statement}{A.66}{Eventually divisible valuation profiles cannot support ghosts}{Proved}
Let $E/\Q$ be any elliptic curve, let $N\ge2$, put
$U=E-E[N]$, and let $R\in E(\Q)$ have infinite order.
Suppose that $a=(a_p)_p\in\widehat\Z$ has only
finitely many zero components and that, for some
prime $\lambda$,
\[
 v_p(a_p)\equiv0\pmod\lambda
       \quad\text{for all but finitely many }p
       \text{ with }a_p\ne0.
 \tag{A66.1}
\]
Then $a\kappa(R)$ has a geometrically abelian Selmer
lift on $U$ if and only if $a$ is a nonzero integer.
Its abelian fibre is the unique point section of $aR$.
The same exact exponent set holds for the intrinsic
metabelian quotient, without a classification of
its fibres.

There is no restriction $\lambda\nmid N$.
In particular this excludes every nonintegral
componentwise $\lambda$th power $a=b^\lambda$
whose zero-component set is finite, and the
concrete exponents
\[
 a_p=p^r\quad\text{for every coefficient prime }p,
                   \qquad r\ge2.
 \tag{A66.2}
\]
The argument uses neither CM nor a finite-Sha
hypothesis. Over any number field, its new
contradiction applies when infinitely many of
the finite valuations $v_p(a_p)$ are positive.
\end{statement}
\begin{proof}
We prove the assertion with infinitely many
positive finite valuations over an arbitrary
number field $k$. Pass to one finite Galois
extension $L/k$ splitting $E[N]$. Choose
independent $P,Q\in E[N](L)$. The four points
$O,P,Q,P+Q$ are distinct, even when $N=2$.
There is a rational function $g\in L(E)^\times$
with
\[
 \operatorname{div}(g)=(P)+(Q)-(P+Q)-(O).
 \tag{A66.3}
\]
Indeed this divisor has degree zero and zero
Abel--Jacobi class. In a Weierstrass plane model
one may take the line through $P,Q$ divided by
the vertical line through $P+Q$: the third
intersection point $-(P+Q)$ cancels.
For each $d\in E[N](L)$ choose its translate
$g_d$ with
\[
 \operatorname{div}(g_d)
   =(P+d)+(Q+d)-(P+Q+d)-(d).
 \tag{A66.4}
\]
These are units on $U$, and every cusp $d$
is a simple pole of the corresponding $g_d$.
The use of simple orders is what permits
every coefficient prime $\lambda$, including
the primes dividing $N$.

A putative global geometrically abelian
Selmer lift gives a class in
$H^1(L,\mu_\lambda)$ for each $g_d$.
Choose representatives of these finitely many
classes. Let $S$ be one finite set of places
of $L$ containing their supports, the places
of bad reduction, the places above $2,3,N$,
the rational primes ramified in $L$, and the
vertical divisor components of all $g_d$.
Also include the finitely many zero-component
primes of $a$ and the finitely many exceptions
to (A66.1). Outside $S$, the points of $E[N]$
are disjoint sections of a smooth model,
and (A66.4) is its exact horizontal divisor.

At $v\notin S$ above the rational prime $q$,
the full compact local Kummer map is injective
on $E(L_v)$. Thus it identifies the forced
local centre of the Selmer lift with
\[
 Q_v=[a]R.
 \tag{A66.5}
\]
This point is nontorsion. The elliptic
logarithm on the compact local group has
torsion kernel; it is nonzero on $R$ and
its value at $[a]R$ is multiplied by the
nonzero scalar $a_q$. Consequently the
local Selmer class is an interior point
section, and the ordinary values of the
units satisfy
\[
 v(g_d(Q_v))\equiv0\pmod\lambda
                   \qquad(d\in E[N](L)).
 \tag{A66.6}
\]
If $Q_v$ reduces to a cusp $c$, the function
$g_c$ has a simple pole there. Its valuation
is the negative of the positive formal
distance from $Q_v$ to $c$. Equation
(A66.6) therefore makes that distance
divisible by $\lambda$.

There are infinitely many distinct ordinary
positive integers $n_j$ which are
$\lambda$th powers and divide $a$ in
$\widehat\Z$. To construct them explicitly,
enumerate distinct primes $p_i$ outside
the finite exceptional set with
$v_{p_i}(a_{p_i})>0$. By (A66.1) these
valuations are at least $\lambda$, and
we may take
\[
 n_j=(p_1\cdots p_j)^\lambda,\qquad
 b_j=a/n_j\in\widehat\Z,\qquad R_j=[n_j]R.
 \tag{A66.7}
\]
Division here is genuine componentwise
division in $\widehat\Z$: at selected
primes the valuation of $n_j$ is
$\lambda\le v_p(a_p)$; at other primes
it is zero. A zero component of $a$
is divisible by every $n_j$, so causes
no obstruction to the definition.
It was placed in $S$ solely to keep the
local centres outside $S$ nontorsion.

Suppose $R_j$ reduces to $d\in E[N](L)$
at $v\notin S$. In the actual compact
local group,
\[
 Q_v-[b_j]d=[b_j](R_j-d),\qquad
 [b_j]d=[b_j\bmod N]d\in E[N](L).
 \tag{A66.8}
\]
In particular $Q_v$ reduces to the cusp
$[b_j]d$. On the formal group the action
of $b_j$ factors through $b_{j,q}\in\Z_q$.
The field $L_v$ is unramified over
$\Q_q$ and $q\ge5$, by the choice of
$S$. The formal logarithm and its inverse
preserve valuation on the formal kernel.
Therefore multiplication by $b_j$ changes
the positive formal distance by exactly
\[
 v_q(b_{j,q})
    =v_q(a_q)-v_q(n_j)
    \in\lambda\Z_{\ge0}.
 \tag{A66.9}
\]
The scalar is nonzero, so both distances
are finite. The distance at $Q_v$ is
divisible by $\lambda$ by (A66.6);
(A66.8)--(A66.9) imply that the distance
of $R_j$ from $d$ is divisible by
$\lambda$ also.

Fix the single function $g$ of (A66.3).
At every $v\notin S$, either $R_j$
does not reduce to its divisor, giving
valuation zero, or its valuation is
plus or minus a cusp distance just
proved divisible by $\lambda$. Hence
\[
 v(g(R_j))\equiv0\pmod\lambda
           \quad\text{for every }j
           \text{ and every }v\notin S.
 \tag{A66.10}
\]
The group
\[
 L(S,\lambda)=
 \{c\in L^\times/(L^\times)^\lambda:
          v(c)\equiv0\pmod\lambda
          \text{ for all }v\notin S\}
 \tag{A66.11}
\]
is finite. Its unit part is finite by
the $S$-unit theorem modulo
$\lambda$th powers; its ideal-class
part lies in the $\lambda$-torsion
of the finite $S$-class group.

By (A66.10), the infinitely many
distinct points $R_j$ lift to the
finite collection of twists
\[
 C_c:\quad z^\lambda=c^{-1}g,
                    \qquad c\in L(S,\lambda).
 \tag{A66.12}
\]
Each smooth projective normalization
is geometrically connected, since $g$
has simple zeros. Its degree-$\lambda$
map to $E$ is fully ramified at the
four distinct points in (A66.3) and
nowhere else. Riemann--Hurwitz gives
\[
 2g(C_c)-2=4(\lambda-1),
                  \qquad g(C_c)=2\lambda-1\ge3.
 \tag{A66.13}
\]
Faltings's theorem makes the rational
points on each member finite
\citep[Satz~7, p.~365, and the correction]{AMFaltings1983,AMFaltings1984Erratum},
contradicting the infinitely many
$R_j$. This proves the new exclusion.

Over $\Q$, if there are only finitely
many positive finite valuations,
the assumed finite zero-component set
makes $a$ a unit at almost all primes.
A.65 applies and forces $a$ to be a
nonzero integer. Conversely every
such integer gives the actual point
section at $aR$, and A.64 supplies
the unique abelian fibre. Projection
to the abelian quotient and actual
point-section existence give the
stated metabelian exponent set.
For (A66.2), choose any prime
$\lambda$ dividing $r$.
\end{proof}

This is an obstruction with genuinely
infinite nonunit support. It does not
replace the fixed finite set $S$ by an
infinite set, and it does not infer
finite support for a general completed
multiplicative class. Among exponents
with finitely many zero components,
an undecided nonintegral exponent must
have, for \emph{every} prime $\lambda$,
infinitely many finite valuations not
divisible by $\lambda$. The elementary
profile $v_p(a_p)=1$ for every $p$ is
not excluded by A.66. Exponents with
infinitely many zero components also
require a different argument: at those
residue primes the compact local
centre can be torsion, and the finite
formal-distance comparison in
(A66.9) ceases to apply.
